\documentclass[tikz,border=8pt]{standalone}
\usepackage{tikz}
\usepackage{amsmath}
\usepackage{amssymb}
\usepackage{pgfplots}
\pgfplotsset{compat=1.18}
\usepackage{ctex}
\usetikzlibrary{calc,positioning,arrows.meta,matrix}

% LC 739 每日温度（单调栈）
% 温度数组 [73,74,75,71,69,72,76,73]
% 单调递减栈存索引，弹出时计算等待天数
% 4 个快照展示栈演化

\begin{document}
\begin{tikzpicture}[
  every node/.style={font=\small},
  tempcell/.style={draw=gray!60,rectangle,minimum width=0.82cm,
    minimum height=0.65cm,thick,fill=blue!5},
  currcell/.style={draw=orange!70,rectangle,minimum width=0.82cm,
    minimum height=0.65cm,very thick,fill=orange!12},
  rescell/.style={draw=green!60!black,rectangle,minimum width=0.82cm,
    minimum height=0.65cm,thick,fill=green!8},
  stackcell/.style={draw=gray!70,rectangle,minimum width=1.1cm,
    minimum height=0.6cm,thick,fill=blue!8},
  popcell/.style={draw=red!60!black,rectangle,minimum width=1.1cm,
    minimum height=0.6cm,very thick,fill=red!8},
  ptr/.style={->,>=Stealth,thick},
]

\def\dx{0.94} % cell width spacing for temperature array

%% ── 标题 ──────────────────────────────────────────────────
\node[font=\large\bfseries] at (5.5, 2.8)
  {每日温度（LC 739 Daily Temperatures）};
\node[font=\small,gray] at (5.5, 2.1)
  {单调递减栈（存索引），弹出时记录等待天数；红箭头连被弹出元素与当前元素};

%% ── 温度数组（顶部参考） ─────────────────────────────────
\node[font=\scriptsize,gray,anchor=east] at (0.6, 1.2) {temp:};
\foreach \i/\v in {0/73,1/74,2/75,3/71,4/69,5/72,6/76,7/73} {
  \node[tempcell] (T\i) at (1.0+\i*\dx, 1.2) {\v};
  \node[font=\tiny,gray] at (1.0+\i*\dx, 1.62) {\i};
}

%% ── 结果数组（顶部参考，紧接温度数组下方） ─────────────────────────
\node[font=\scriptsize,gray,anchor=east] at (0.6, 0.3) {ans:};
\foreach \i/\v in {0/1,1/1,2/4,3/2,4/1,5/1,6/0,7/0} {
  \node[rescell] (R\i) at (1.0+\i*\dx, 0.3) {\v};
}

%% ── 快照说明区域（4 个水平排列的快照） ──────────────────
%% 每个快照：左侧是栈（竖排，右为栈顶），右侧是说明
%% 快照 A：处理 i=0..2，栈 [0,1,2]（存索引，温度递增不弹出）

\def\snapYtop{-2.2}
\def\snapYstep{-3.5}

%% ══════════════════════════════════════════════════════════
%% 快照 1：处理 i=0(73) i=1(74) i=2(75)，均压栈
%% 栈（底→顶）：[0,1,2]，无弹出
%% ══════════════════════════════════════════════════════════
\def\ya{\snapYtop}
\node[font=\scriptsize\bfseries,blue!70] at (0.4, \ya-0.1)
  {快照1: i=0→2 全部压栈};

% 当前处理位置高亮
\node[currcell] at (1.0+2*\dx, 1.2) {75};

% 栈（横排，左=底，右=顶）
\node[font=\tiny,gray,anchor=east] at (0.6,\ya-0.8) {栈:};
\node[stackcell,font=\scriptsize] (Sa0) at (1.1, \ya-0.8) {0(73)};
\node[stackcell,font=\scriptsize] (Sa1) at (2.3, \ya-0.8) {1(74)};
\node[stackcell,font=\scriptsize] (Sa2) at (3.5, \ya-0.8) {2(75)};
\node[font=\tiny,green!60!black,anchor=west] at (4.1, \ya-0.8) {$\leftarrow$ top};
\draw[thick,gray!60] (0.65, \ya-1.15) -- (4.1, \ya-1.15);
\node[font=\tiny,gray,anchor=west] at (0.65, \ya-1.35) {栈底←};

\node[draw=gray!40,fill=white,thin,rounded corners=2pt,
      font=\scriptsize,inner sep=4pt,align=left]
  at (7.8, \ya-0.7)
  {温度递增，逐个压栈\\无弹出，ans 暂未更新};

%% ══════════════════════════════════════════════════════════
%% 快照 2：处理 i=3(71)，不弹出（71<75），压入 3
%% 随后 i=4(69)，不弹出（69<75），压入 4
%% ══════════════════════════════════════════════════════════
\def\yb{\snapYtop+\snapYstep}
\node[font=\scriptsize\bfseries,blue!70] at (0.4, \yb-0.1)
  {快照2: i=3(71),4(69) 压栈};

\node[currcell] at (1.0+4*\dx, 1.2) {69};

\node[font=\tiny,gray,anchor=east] at (0.6,\yb-0.8) {栈:};
\node[stackcell,font=\scriptsize] (Sb0) at (1.1, \yb-0.8) {0(73)};
\node[stackcell,font=\scriptsize] (Sb1) at (2.3, \yb-0.8) {1(74)};
\node[stackcell,font=\scriptsize] (Sb2) at (3.5, \yb-0.8) {2(75)};
\node[stackcell,font=\scriptsize] (Sb3) at (4.7, \yb-0.8) {3(71)};
\node[stackcell,font=\scriptsize] (Sb4) at (5.9, \yb-0.8) {4(69)};
\node[font=\tiny,green!60!black,anchor=west] at (6.5, \yb-0.8) {$\leftarrow$ top};
\draw[thick,gray!60] (0.65, \yb-1.15) -- (6.5, \yb-1.15);
\node[font=\tiny,gray,anchor=west] at (0.65, \yb-1.35) {栈底←};

\node[draw=gray!40,fill=white,thin,rounded corners=2pt,
      font=\scriptsize,inner sep=4pt,align=left]
  at (9.2, \yb-0.7)
  {71,69 均小于栈顶 75\\不触发弹出，继续压栈};

%% ══════════════════════════════════════════════════════════
%% 快照 3：处理 i=5(72)，弹出 4(69),3(71)
%% ans[4]=5-4=1, ans[3]=5-3=2
%% ══════════════════════════════════════════════════════════
\def\yc{\snapYtop+2*\snapYstep}
\node[font=\scriptsize\bfseries,blue!70] at (0.4, \yc-0.1)
  {快照3: i=5(72)，弹出 4(69) 和 3(71)};

\node[currcell] at (1.0+5*\dx, 1.2) {72};

% 被弹出的格子
\node[font=\tiny,gray,anchor=east] at (0.6,\yc-0.8) {栈:};
\node[stackcell,font=\scriptsize] (Sc0) at (1.1, \yc-0.8) {0(73)};
\node[stackcell,font=\scriptsize] (Sc1) at (2.3, \yc-0.8) {1(74)};
\node[stackcell,font=\scriptsize] (Sc2) at (3.5, \yc-0.8) {2(75)};
\node[stackcell,font=\scriptsize] (Sc5) at (4.7, \yc-0.8) {5(72)};
\node[font=\tiny,green!60!black,anchor=west] at (5.3, \yc-0.8) {$\leftarrow$ top};

% 弹出的两个（显示在栈下方，半透明红色填充表示"已弹出"）
\node[popcell,font=\scriptsize,fill=red!15] (Sc3pop) at (4.7, \yc-1.7) {3(71)};
\node[font=\tiny,red!70!black,anchor=west] at ($(Sc3pop.east)+(0.15,0)$) {弹出 → ans[3]=5-3=\textbf{2}};

\node[popcell,font=\scriptsize,fill=red!15] (Sc4pop) at (4.7, \yc-2.4) {4(69)};
\node[font=\tiny,red!70!black,anchor=west] at ($(Sc4pop.east)+(0.15,0)$) {弹出 → ans[4]=5-4=\textbf{1}};

\draw[thick,gray!60] (0.65, \yc-1.15) -- (5.3, \yc-1.15);
\node[font=\tiny,gray,anchor=west] at (0.65, \yc-1.35) {栈底←};

\node[draw=gray!40,fill=white,thin,rounded corners=2pt,
      font=\scriptsize,inner sep=4pt,align=left]
  at (9.5, \yc-0.9)
  {72 > 69 弹出 4，ans[4]=5-4=1\\72 > 71 弹出 3，ans[3]=5-3=2\\72 < 75 停止弹出，压入 5};

%% ══════════════════════════════════════════════════════════
%% 快照 4：处理 i=6(76)，弹出 5(72),2(75),1(74),0(73)
%% ans[5]=6-5=1, ans[2]=6-2=4, ans[1]=6-1=5, ans[0]=6-0=6 (→但只更 0,1,2,5)
%% ══════════════════════════════════════════════════════════
\def\yd{\snapYtop+3*\snapYstep}
\node[font=\scriptsize\bfseries,blue!70] at (0.4, \yd-0.1)
  {快照4: i=6(76)，弹出所有，压 6；i=7(73) 压栈};

\node[currcell] at (1.0+6*\dx, 1.2) {76};

\node[font=\tiny,gray,anchor=east] at (0.6,\yd-0.8) {栈:};
\node[stackcell,font=\scriptsize] (Sd6) at (1.1, \yd-0.8) {6(76)};
\node[stackcell,font=\scriptsize] (Sd7) at (2.3, \yd-0.8) {7(73)};
\node[font=\tiny,green!60!black,anchor=west] at (2.9, \yd-0.8) {$\leftarrow$ top};
\draw[thick,gray!60] (0.65, \yd-1.15) -- (2.9, \yd-1.15);
\node[font=\tiny,gray,anchor=west] at (0.65, \yd-1.35) {栈底←};

% 快照4：依次弹出 5/2/1/0，对应 ans 已在结果数组中（去除 4 条重叠弧形箭头，避免视觉混乱）
% 弹出顺序在右侧说明框中文字描述

\node[draw=gray!40,fill=white,thin,rounded corners=2pt,
      font=\scriptsize,inner sep=4pt,align=left]
  at (8.8, \yd-0.8)
  {76 大于栈中所有温度\\依次弹出，记录 ans[5..0]\\最终压入 6(76)，再压 7(73)\\7(73) 无更热天，ans[6]=ans[7]=0};

%% ── 分隔线 ──────────────────────────────────────────────
\draw[gray!25,dashed] (0.0, -16.5) -- (13.0, -16.5);

%% ── 算法说明（顶部 ans 数组已展示最终结果，此处仅放算法说明）
\node[draw=gray!50,fill=white,rounded corners=3pt,font=\small,
      inner sep=8pt,align=left]
  at (6.5, -17.8)
  {\textbf{单调递减栈（O(n) 时间，O(n) 空间）：}\\[4pt]
   维护单调递减栈（存索引 i），遍历 temp 数组：\\[2pt]
   \quad 当 temp[i] > temp[stack.top()] 时循环弹出 j=stack.pop()，\\
   \quad 记录 ans[j] = i - j（等待天数）\\[2pt]
   \quad 将 i 压栈；遍历结束，栈中剩余索引对应 ans=0};

\end{tikzpicture}
\end{document}
